\section{Usage}
\label{sec:usage}

The API has two levels: a bare loop over \code{model.update} for people who want control, and
a \clsname{Trainer} for people who want Keras. Both are shown below; every listing is
runnable as written, and all of them have full counterparts in the repository's
\code{examples/} directory.

\subsection{Thirty seconds}

\begin{lstlisting}[style=ttpython]
import torchtsetlin as tt

x_train, y_train = tt.data.make_noisy_xor(5000, noise=0.4)    # Boolean features, int labels
x_test,  y_test  = tt.data.make_noisy_xor(2000, noise=0.0, seed=1)

model = tt.TsetlinMachine(n_features=12, n_classes=2, n_clauses=20, T=15, s=3.9).to("cuda")

# update() is the analogue of loss.backward(); optimizer.step()
for epoch in range(50):
    model.train()
    for i in range(0, 5000, 10):
        model.update(x_train[i:i+10].cuda(), y_train[i:i+10].cuda())

model.eval()                                  # empty clauses now vote for nothing
acc = (model(x_test.cuda()).argmax(1) == y_test.cuda()).float().mean()
print(model.rules()[:2])
# ['IF x0 AND NOT x1 THEN 1', 'IF x1 AND NOT x0 THEN 1']
\end{lstlisting}

The last line is the point of the exercise: the model has recovered the generating rule of
the Noisy XOR benchmark from data with 40\% label noise, and printed it.

\subsection{Tabular data, rules in the original units}

Continuous columns become Boolean through a thermometer encoder. Passing the encoder's
feature names to the model is what turns \code{x37} back into \code{worst perimeter>=105.9},
and \code{max\_included\_literals} keeps the printed rules short enough to read.

\begin{lstlisting}[style=ttpython]
from sklearn.datasets import load_breast_cancer
from sklearn.model_selection import train_test_split
import torch, torchtsetlin as tt

tt.seed_everything(0)
data = load_breast_cancer()
x_tr, x_te, y_tr, y_te = train_test_split(data.data, data.target, test_size=0.3, random_state=0)
x_tr, x_te = torch.tensor(x_tr).float(), torch.tensor(x_te).float()

enc = tt.data.ThermometerEncoder(n_bits=8, strategy="quantile").fit(x_tr)
xb_tr, xb_te = enc(x_tr), enc(x_te)

model = tt.TsetlinMachine(xb_tr.shape[1], 2, n_clauses=100, T=20, s=4.0,
                          max_included_literals=6)
model.feature_names = enc.feature_names(list(data.feature_names))   # readable literals
model.class_names   = list(data.target_names)

trainer = tt.Trainer(model, batch_size=8, callbacks=[tt.EarlyStopping(patience=30)])
trainer.fit((xb_tr, y_tr), epochs=200, val_data=(xb_te, y_te))
print(trainer.test((xb_te, y_te)))                 # {'test_accuracy': 0.96...}

for rule in model.rules(polarity=1)[:5]:           # clauses voting *for* a class
    print(rule)

# Which original columns carry the model, with the thermometer bits summed back together
imp = tt.interpret.global_feature_importance(model, class_index=0)
per_feature = tt.interpret.aggregate_literal_importance(
    imp, [i // 8 for i in range(xb_tr.shape[1])])
\end{lstlisting}

\subsection{Convolutional MNIST on a GPU}

The convolutional path differs in three places only: the encoder (adaptive thresholding
rather than a fixed cut), the model class, and the patch geometry. Keeping the encoder on the
compute device means the images never travel back to the host.

\begin{lstlisting}[style=ttpython]
import torch, torchtsetlin as tt
dev = "cuda"

enc = tt.data.AdaptiveThresholdEncoder(block_size=11, C=2, method="gaussian").to(dev)
booleanize = lambda x: enc(x * 255.0)

x_train, y_train = tt.data.load_mnist_boolean("./data", train=True,
                                              encoder=booleanize, device=dev)
x_test,  y_test  = tt.data.load_mnist_boolean("./data", train=False,
                                              encoder=booleanize, device=dev)

model = tt.ConvTsetlinMachine(
    n_classes=10, n_clauses=200, T=100, s=10.0,
    patch_size=10, stride=1, position_encoding=True, weighted=True,
    input_shape=tuple(x_train.shape[1:]),           # (1, 28, 28)
)

trainer = tt.Trainer(model, device=dev, batch_size=32, callbacks=[
    tt.EarlyStopping(monitor="val_accuracy", patience=5),
    tt.ModelCheckpoint("mnist_conv.pt", monitor="val_accuracy"),
    tt.CSVLogger("mnist_conv.csv"),
    tt.StateSummaryLogger(),                        # mean literals/clause, empty clauses
])
trainer.fit((x_train, y_train), epochs=10, val_data=(x_test, y_test))

votes = trainer.predict((x_test, y_test), return_votes=True)
print(tt.metrics.classification_report(votes, y_test))

tt.viz.plot_conv_clauses(model, range(16))          # what each clause looks for
\end{lstlisting}

\begin{ttcallout}[Hyper-parameter scheduling]
Tsetlin hyper-parameters are ordinary Python attributes, so annealing one is a callback and
not a framework feature:
\code{tt.HyperparameterSchedule("s", lambda epoch: 10.0 - 5.0 * epoch / 30)}, or the same for
\code{drop\_clause\_p}. \code{model.state\_summary()} reports mean included literals per
clause and the number of empty clauses --- the two numbers that tell you whether $s$ is in the
right range.
\end{ttcallout}

\subsection{Explaining a single prediction}

\begin{lstlisting}[style=ttpython]
expl = tt.interpret.explain(model, x_test[0], feature_names=model.feature_names)
print(expl)                     # shape of the output; the numbers depend on the run
# prediction: 1  votes: [-14.0, 17.0]
#   [+3 -> output 1] clause 42: IF worst perimeter>=105.9 AND NOT mean texture>=21.8
#   [+2 -> output 1] clause 77: IF worst concave points>=0.15
#   ...

torch.save(model.state_dict(), "model.pt")           # ordinary PyTorch checkpointing
reloaded = tt.TsetlinMachine(None, 2, n_clauses=100, T=20, s=4.0)   # lazy: shape from ckpt
reloaded.load_state_dict(torch.load("model.pt"))
\end{lstlisting}

The explanation is not a surrogate: those clauses are the model, and the listed votes add up
to the reported vote sums.

\subsection{Building a new variant}

A new model supplies an output layer and a feedback policy; the base class provides
everything else, including chunking, dropout, the batched commit and the whole workflow
layer.

\begin{lstlisting}[style=ttpython]
class MyTsetlinMachine(tt.TsetlinMachineBase):
    def __init__(self, n_features, n_clauses, T, **kw):
        self.T = T
        super().__init__(n_features, n_clauses, **kw)

    def _votes(self, clause_out, clamp=True):        # (B, C) bool -> (B, n_outputs)
        ...

    def _coerce_targets(self, y, n, device):         # validate / convert targets
        ...

    def _select_feedback(self, votes, y, clause_out):  # -> (type_i, type_ii, aux) masks
        ...
\end{lstlisting}

Or skip the class hierarchy entirely and compose the functional layer directly --- useful when
prototyping a different feedback rule:

\begin{lstlisting}[style=ttpython]
from torchtsetlin import functional as F

lit     = F.to_literals(x)                                  # (B, 2F)
include = (state >= N).float()
c       = F.clause_outputs(lit, include, empty_value=True)  # (B, C)
v       = F.vote_sums(c, weights, T)
n_true, n_false, n_ib = F.type_i_counts(lit, sel_fire_i, sel_nofire_i)
n2      = F.type_ii_counts(lit, sel_fire_ii)
F.apply_feedback(state, n_true, n_false, n_ib, n2, n_states=N, s=s)
\end{lstlisting}

It is plain tensor code, so \code{torch.compile}, a custom CUDA kernel or an alternative
sampling scheme can be dropped in without touching anything else.
